\documentclass[12pt, a4paper]{report}

%% --- Paquetes ---
\usepackage[utf8]{inputenc}
\usepackage[T1]{fontenc}
\usepackage[spanish]{babel}
\usepackage{geometry}
\usepackage{booktabs}
\usepackage{array}
\usepackage{tabularx}
\usepackage{setspace}
\usepackage{parskip}
\usepackage{titlesec}
\usepackage{fancyhdr}
\usepackage{hyperref}
\usepackage{listings}
\usepackage{xcolor}
\usepackage{float}
\usepackage{enumitem}
\usepackage{longtable}
\usepackage{mdframed}
\usepackage{graphicx}

\geometry{top=2.5cm, bottom=2.5cm, left=3cm, right=2.5cm}

\definecolor{grisclaro}{RGB}{245,245,245}
\definecolor{griscode}{RGB}{230,230,230}

\lstset{
  basicstyle=\ttfamily\small,
  backgroundcolor=\color{grisclaro},
  frame=single,
  breaklines=true,
  showstringspaces=false,
  captionpos=b,
  literate={á}{{\'a}}1 {é}{{\'e}}1 {í}{{\'i}}1
           {ó}{{\'o}}1 {ú}{{\'u}}1 {ñ}{{\~n}}1
           {Á}{{\'A}}1 {É}{{\'E}}1 {Í}{{\'I}}1
           {Ó}{{\'O}}1 {Ú}{{\'U}}1 {Ñ}{{\~N}}1
}

\pagestyle{plain}

\titleformat{\chapter}[display]
  {\normalfont\Large\bfseries}
  {\chaptertitlename\ \thechapter}{10pt}{\Huge}
\titlespacing*{\chapter}{0pt}{-10pt}{20pt}

\titleformat{\section}{\large\bfseries}{\thesection}{1em}{}
\titleformat{\subsection}{\normalsize\bfseries}{\thesubsection}{1em}{}

\hypersetup{colorlinks=false, hidelinks}

\begin{document}
\onehalfspacing

\begin{titlepage}
  \thispagestyle{empty}
  \centering
  \vspace*{0.5cm}

  {\Large\bfseries Universidad Técnica Estatal de Quevedo\\[0.4cm]}
  {\normalsize Facultad de Ciencias de la Computación y Diseño Gráfico\\[0.1cm]
   Carrera de Ingeniería en Software\\[1.5cm]}

  \rule{\linewidth}{1.5pt}\\[0.4cm]
  {\LARGE\bfseries
    Sistema de Gestión de Solicitudes\\[0.2cm]
    para Soporte Técnico de Internet}\\[0.4cm]
  \rule{\linewidth}{1.5pt}\\[1cm]

  {\large\bfseries Entrega 2 --- Diseño de Arquitectura de Microservicios\\[0.3cm]}
  {\normalsize Versión 2.0\\[1.5cm]}

  \begin{flushleft}
    \textbf{Equipo de trabajo:}\\[0.4cm]
    \hspace{1cm}
    \begin{minipage}{0.85\textwidth}
      \begin{itemize}
        \setlength{\itemsep}{0.3cm}
        \item Pacheco Cárdenas Cristhian Daniel \quad --- \quad Arquitecto
        \item Carpio Mendoza Carlos José \quad --- \quad Líder de Desarrollo
        \item Álvarez Párraga Jeremy Alexis \quad --- \quad Responsable de Calidad
        \item Cando Moreno Robinson Rodrigo \quad --- \quad Responsable de Documentación
      \end{itemize}
    \end{minipage}
  \end{flushleft}

  \vspace{1cm}

  \begin{flushleft}
    \hspace{1cm}
    \begin{tabular}{ll}
      \textbf{Docente:}            & Gleiston Cicerón Guerrero Ulloa \\[0.3cm]
      \textbf{Asignatura:}         & Aplicaciones Distribuidas \\[0.3cm]
      \textbf{Período Académico:}  & Séptimo Semestre \\[0.3cm]
      \textbf{Fecha de Entrega:}   & Junio 2026 \\
    \end{tabular}
  \end{flushleft}

  \vfill
  {\small Quevedo -- Ecuador \quad | \quad 2026}
\end{titlepage}

\tableofcontents
\newpage

\chapter*{Entrega 1: Propuesta Inicial del Proyecto}
\addcontentsline{toc}{chapter}{Entrega 1: Propuesta Inicial del Proyecto}

\section*{Resumen Ejecutivo}

\subsection*{Problema que se resuelve}

Muchas organizaciones, especialmente de tamaño pequeño y mediano, gestionan
sus solicitudes de soporte técnico a través de canales informales como correos
electrónicos, llamadas telefónicas o mensajería instantánea. Esta forma de
trabajar genera desorganización, pérdida de información y tiempos de respuesta
prolongados. El presente proyecto propone un \textbf{Sistema de Gestión de
Solicitudes para Soporte Técnico de Internet}, orientado a formalizar y
estructurar el proceso de atención de incidencias relacionadas con la
conectividad a Internet dentro de entornos organizacionales.

\vspace{0.5cm}

\subsection*{Usuarios del sistema}

El sistema está dirigido a organizaciones que ofrecen o dependen de servicios
de Internet y que necesitan un mecanismo ordenado para gestionar incidencias
técnicas. Lo utilizarán tanto el cliente que reporta problemas como el equipo
técnico encargado de resolverlos y el personal administrativo que supervisa y
coordina el proceso de atención.

\vspace{0.5cm}

\subsection*{Un sistema distribuido es la solución apropiada}

Si esta aplicación se construyera como un sistema único y centralizado,
todo dependería de un solo punto: si ese punto falla, el servicio deja de
funcionar por completo.

Un enfoque distribuido resuelve esto dividiendo la aplicación en partes
separadas: la interfaz que ve el usuario, la lógica que procesa las
solicitudes y la base de datos donde se guarda la información. Cada parte
trabaja de forma independiente y se comunican entre sí. Esto trae beneficios
directos: si una parte falla, las demás pueden seguir funcionando; si la
cantidad de usuarios crece, se puede reforzar solo la parte que lo necesita;
y si se requiere hacer un cambio, se modifica únicamente la parte
correspondiente sin afectar el resto del sistema.

\newpage

\subsection*{Alcance del período académico}

Durante el período académico, el trabajo consistirá en tomar la aplicación web del Sistema de Gestión de Solicitudes para Soporte Técnico de Internet, ya existente, e identificar de qué manera puede reorganizarse bajo un enfoque distribuido. A partir de ese análisis, se irán aplicando gradualmente los conceptos trabajados a lo largo del semestre sobre su estructura y diseño: cómo se comunican sus componentes, cómo se organiza y distribuye su información, y de qué forma puede estructurarse en capas independientes.

La idea es transformar la aplicación web original en un ecosistema basado en microservicios desacoplados y orientados a eventos, dividiendo la plataforma en componentes individuales de autenticación, gestión de reportes, notificaciones e inteligencia artificial que interactuarán de forma que no esperarse el uno al otro para seguir trabajando.

\newpage

\section{Mejora de decisiones tomadas}

\begin{longtable}{|p{2.8cm}|p{4.2cm}|p{5.8cm}|}
\hline
\textbf{Aspecto} & \textbf{Entrega 1} & \textbf{Entrega 2} \\
\hline
\endfirsthead
\hline
\textbf{Aspecto} & \textbf{Entrega 1} & \textbf{Entrega 2} \\
\hline
\endhead
Número de servicios &
Mencionados de forma general, sin número exacto &
Se definen 5 servicios con nombre, puerto y
responsabilidad clara cada uno \\
\hline
Bases de datos &
No se mencionaba nada &
Cada servicio gestiona su propia base de datos.
Se usa PostgreSQL para datos estructurados y
MongoDB donde la información es más flexible \\
\hline
Comunicación &
No estaba definida &
REST para operaciones que necesitan respuesta
inmediata; Kafka para avisos en segundo plano
que no requieren esperar \\
\hline
Categorización &
Mencionada como idea general &
Se define un servicio dedicado que sugiere la
categoría de una solicitud nueva basándose en
las categorías ya registradas en el sistema \\
\hline
Punto de entrada &
No definido &
Se incorpora un API Gateway como puerta de
entrada única para todos los clientes \\
\hline
Despliegue &
No mencionado &
Se define una estrategia con Docker para que
cada servicio pueda ejecutarse de forma
independiente en cualquier máquina \\
\hline
\end{longtable}

\newpage

\chapter{Arquitectura de Microservicios}

\section{¿Qué es un microservicio?}

Un microservicio es una parte pequeña y autónoma de un sistema más
grande. En lugar de tener una sola aplicación que lo haga todo,
el trabajo se divide en piezas. Cada pieza tiene su propia función,
su propio código y su propia base de datos. Si una pieza falla, las
demás siguen funcionando.

\section{¿Por qué se eligió esta arquitectura para el proyecto?}

La aplicación original del sistema de soporte técnico funciona como un
bloque único. Eso tiene un problema, debido a que si hay un error en cualquier
parte, todo el sistema puede quedar fuera de servicio.

Se eligió una arquitectura de microservicios por tres razones concretas:

\begin{enumerate}
  \item \textbf{Separación de responsabilidades.} Cada parte del
        sistema hace una sola cosa y la hace bien. El servicio de
        autenticación no sabe nada de los tickets, y el servicio de
        notificaciones no sabe nada de los usuarios.

  \item \textbf{Independencia ante fallos.} Si el servicio de
        notificaciones tiene un problema, los clientes aún pueden
        crear solicitudes y los técnicos pueden atenderlas. El fallo
        queda contenido.

  \item \textbf{Facilidad para crecer.} Si en el futuro hay más
        solicitudes de las que el sistema puede manejar, se puede
        reforzar únicamente el servicio de solicitudes sin tocar
        el resto.
\end{enumerate}

\newpage
\section{Los cinco servicios del sistema}

Después del análisis del equipo, se definieron cinco servicios.
Cada uno corresponde a una responsabilidad distinta del sistema.

\begin{longtable}{|p{3.8cm}|p{2cm}|p{2.2cm}|p{4.5cm}|}
\hline
\textbf{Servicio} & \textbf{Puerto} & \textbf{Base de datos} &
\textbf{¿Qué hace?} \\
\hline
\endfirsthead
\hline
\textbf{Servicio} & \textbf{Puerto} & \textbf{Base de datos} &
\textbf{¿Qué hace?} \\
\hline
\endhead
auth-service & 8001 & PostgreSQL &
Registra usuarios, verifica contraseñas y
genera el token de acceso \\
\hline
ticket-service & 8002 & PostgreSQL &
Crea, actualiza y cierra solicitudes de soporte.
Es el servicio central del sistema \\
\hline
notification-service & 8003 & MongoDB &
Envía avisos al cliente y al técnico cuando
algo cambia en la solicitud \\
\hline
ai-service & 8004 & MongoDB &
Sugiere automáticamente la categoría de una
solicitud nueva comparando su descripción
con las categorías existentes \\
\hline
report-service & 8005 & PostgreSQL &
Muestra estadísticas: solicitudes atendidas,
tiempos de respuesta y rendimiento por técnico \\
\hline
\end{longtable}


\section{Cada servicio tiene su propia base de datos}

Si todos los servicios compartieran una sola base de datos, estarían
conectados entre sí de una forma invisible pero peligrosa y si ocurre un cambio
en la estructura de la base de datos podría romper varios servicios
al mismo tiempo, aunque ese cambio fuera pequeño.

Al dar a cada servicio su propia base de datos, cada uno puede
evolucionar de forma independiente. El equipo puede cambiar cómo
guarda las solicitudes sin afectar en absoluto al servicio de
notificaciones ni al de reportes.

\section{Diagrama de la arquitectura}

El siguiente diagrama muestra cómo se conectan todos los componentes
del sistema. Los clientes siempre entran por el API Gateway. Los
servicios que necesitan comunicarse entre sí lo hacen a través de
Kafka cuando no requieren respuesta inmediata.
\begin{figure}[H]
\centering
\includegraphics[width=\textwidth]{diagrama_arquitectura.png}
\caption{Diagrama de arquitectura del sistema distribuido}
\end{figure}

\newpage
\section{La comunicación entre lo servicios}

Los servicios se comunican de dos formas distintas, despendiendo de si la
operación necesita una respuesta inmediata o no.

\subsection{Comunicación directa (REST)}

Cuando el cliente necesita una respuesta para continuar, se usa REST.
Por ejemplo, para iniciar sesión el sistema debe confirmar de inmediato
si las credenciales son correctas. No puede esperar. Lo mismo ocurre
cuando se crea una solicitud: el cliente necesita recibir el número
de su caso de forma instantánea.

Las rutas que usan REST son las que van del API Gateway hacia
auth-service, ticket-service y report-service.

\subsection{Comunicación por eventos (Kafka)}

Cuando una operación puede resolverse en segundo plano sin que el
cliente espere, se usa Apache Kafka. Kafka funciona como un tablón
de anuncios: un servicio publica un aviso y otros servicios lo
recogen cuando pueden.

Por ejemplo, cuando se crea una solicitud nueva:

\begin{enumerate}
  \item El ticket-service la registra y responde al cliente de
        inmediato con el número de caso.
  \item En paralelo, publica un aviso en Kafka diciendo
        ``nueva solicitud creada''.
  \item El ai-service recibe ese aviso y sugiere una categoría
        en segundo plano.
  \item El notification-service también lo recibe y le envía
        un correo de confirmación al cliente.
\end{enumerate}

El cliente no espera ninguno de esos pasos adicionales. Ya tiene
su número de caso y el resto ocurre de forma automática.

\newpage
\section{Los estados de una solicitud}

Una solicitud pasa por cuatro estados a lo largo de su vida:

\begin{enumerate}
  \item \textbf{Pendiente:} La solicitud fue creada pero aún no
        se le asignó ningún técnico.
  \item \textbf{En Proceso:} Un técnico está trabajando en el caso.
  \item \textbf{Resuelta:} El técnico registró la solución.
  \item \textbf{Cerrada:} El caso quedó formalmente cerrado.
\end{enumerate}

Cada cambio de estado queda registrado automáticamente en el
historial de la solicitud, junto con la fecha y el usuario que
realizó el cambio.

\newpage

\chapter{Base de Datos del Sistema}


\section{Estructura general}

La base de datos original del sistema fue diseñada para una aplicación
única. Al migrar a microservicios, esa estructura sirve como punto de
partida para definir qué información pertenece a cada servicio.

A continuación se describen las tablas más importantes del sistema,
organizadas según el servicio al que corresponden.

\section{Tablas del auth-service}

Este servicio gestiona todo lo relacionado con identidad y acceso.
Sus tablas principales son: usuario y rol.

La tabla usuario guarda el nombre, correo, contraseña
(cifrada) y el rol de cada persona. La tabla rol define
los cuatro tipos de usuario del sistema: Superusuario, Administrador,
Técnico y Cliente.

Adicionalmente existen las tablas cliente y técnico,
que extienden la información básica del usuario. La tabla cliente
guarda la dirección del cliente. La tabla técnico guarda la
especialidad, el nivel de soporte (L1, L2 o L3) y si está habilitado
para recibir casos.

\begin{longtable}{|p{3.5cm}|p{4cm}|p{5cm}|}
\hline
\textbf{Tabla} & \textbf{Campos principales} & \textbf{¿Para qué sirve?} \\
\hline
\endfirsthead
\hline
\textbf{Tabla} & \textbf{Campos principales} & \textbf{¿Para qué sirve?} \\
\hline
\endhead
usuario &
id, nombre, correo, contraseña, id\_rol &
Centraliza la identidad de todos los usuarios \\
\hline
rol &
id, nombre\_rol &
Define los perfiles de acceso al sistema \\
\hline
cliente &
id\_usuario, dirección &
Extiende al usuario con información del cliente \\
\hline
técnico &
id\_usuario, especialidad, nivel, habilitado &
Extiende al usuario con información del técnico \\
\hline
\end{longtable}

\section{Tablas del ticket-service}

Este es el servicio central del sistema. Gestiona el ciclo de vida
completo de cada solicitud de soporte.

La tabla solicitud es el corazón: guarda la descripción del
problema, la categoría, el cliente que la creó, la fecha y el estado
actual. La tabla estado define los cuatro estados posibles.
La tabla categoría contiene los tipos de problema que el
sistema puede atender: Sin conexión a Internet, Baja velocidad
y Cortes intermitentes.

El historial de cambios de estado se guarda en la tabla
\textbf{historial\_estado}, que registra automáticamente cada
transición mediante un trigger en la base de datos.

La tabla \textbf{asignacion\_solicitud} registra a qué técnico
o grupo de técnicos se le asignó cada caso. La tabla
\textbf{reporte\_solicitud} guarda el detalle de lo que el técnico
hizo para resolver el problema.

\begin{longtable}{|p{4cm}|p{4cm}|p{4.5cm}|}
\hline
\textbf{Tabla} & \textbf{Campos principales} & \textbf{¿Para qué sirve?} \\
\hline
\endfirsthead
\hline
\textbf{Tabla} & \textbf{Campos principales} & \textbf{¿Para qué sirve?} \\
\hline
\endhead
solicitud &
id, descripción, id\_categoría, id\_cliente, id\_estado\_actual &
Registro principal de cada caso de soporte \\
\hline
estado &
id, nombre\_estado &
Define los estados: Pendiente, En Proceso, Resuelta, Cerrada \\
\hline
categoría &
id, nombre\_categoría &
Clasifica el tipo de problema reportado \\
\hline
historial\_estado &
id, id\_solicitud, id\_estado, fecha\_hora\_cambio &
Registra cada cambio de estado automáticamente \\
\hline
asignacion\_solicitud &
id, id\_solicitud, id\_técnico, id\_grupo, fecha\_inicio &
Vincula una solicitud con el técnico o grupo asignado \\
\hline
reporte\_solicitud &
id, id\_solicitud, id\_técnico, detalle\_reporte &
Documenta lo que se hizo para resolver el caso \\
\hline
adjunto &
id, id\_solicitud, nombre\_archivo, tipo\_archivo &
Archivos (PNG, JPG, PDF) adjuntados a una solicitud \\
\hline
\end{longtable}

\section{Tablas del report-service}

El servicio de reportes consulta información consolidada del sistema.
Trabaja principalmente con los datos de solicitudes, asignaciones e
historial para construir las métricas que el administrador necesita ver.

En la arquitectura distribuida, este servicio tendrá acceso a una
copia de lectura de los datos relevantes, actualizada a través de
los eventos que publica el ticket-service en Kafka.

\section{Tablas del notification-service}

Este servicio no necesita acceder a las tablas principales del sistema.
Recibe la información que necesita directamente desde los eventos de
Kafka: el nombre del destinatario, el tipo de evento y el mensaje
a enviar. Guarda únicamente el historial de notificaciones enviadas
en su propia base de datos MongoDB.

\newpage

\chapter{Contrato del API REST}

Cuando la aplicación web necesita pedirle algo al servidor por
ejemplo, crear una solicitud o consultar su estado necesitan ponerse
de acuerdo en cómo hablar. Ese acuerdo escrito
se llama contrato del API.

Definirlo antes de programar evita confusiones entre los miembros
del equipo y facilita que el trabajo avance en paralelo: quien
desarrolla la interfaz y quien desarrolla el servidor saben
exactamente qué esperar del otro.

\section{Reglas generales}

Todas las rutas del sistema siguen las mismas convenciones:

\begin{itemize}
  \item Todas empiezan con \texttt{/api/v1/}
  \item Las rutas protegidas requieren un token en el encabezado:\\
        \texttt{Authorization: Bearer \{token\}}
  \item Las respuestas siempre vienen en formato JSON
  \item Se usan los códigos HTTP estándar: 200 (éxito), 201 (creado),
        400 (datos incorrectos), 401 (sin permiso), 404 (no encontrado)
\end{itemize}

\section{Servicio de autenticación}

\begin{longtable}{|p{1.8cm}|p{4.8cm}|p{4.2cm}|p{2.5cm}|}
\hline
\textbf{Método} & \textbf{Ruta} & \textbf{¿Qué hace?} &
\textbf{Acceso} \\
\hline
\endfirsthead
\hline
\textbf{Método} & \textbf{Ruta} & \textbf{¿Qué hace?} &
\textbf{Acceso} \\
\hline
\endhead
POST & /api/v1/auth/register &
Crea una cuenta nueva & Público \\
\hline
POST & /api/v1/auth/login &
Verifica credenciales y entrega el token & Público \\
\hline
POST & /api/v1/auth/logout &
Cancela la sesión activa & Autenticado \\
\hline
GET  & /api/v1/auth/validate &
Verifica si el token es válido & Solo interno \\
\hline
\end{longtable}

\section{Servicio de solicitudes}

\begin{longtable}{|p{1.8cm}|p{5.2cm}|p{3.8cm}|p{1.8cm}|}
\hline
\textbf{Método} & \textbf{Ruta} & \textbf{¿Qué hace?} &
\textbf{Acceso} \\
\hline
\endfirsthead
\hline
\textbf{Método} & \textbf{Ruta} & \textbf{¿Qué hace?} &
\textbf{Acceso} \\
\hline
\endhead
POST  & /api/v1/tickets &
Registra una nueva solicitud & Cliente \\
\hline
GET   & /api/v1/tickets &
Lista solicitudes según el rol del usuario & Todos \\
\hline
GET   & /api/v1/tickets/\{id\} &
Muestra el detalle de una solicitud & Todos \\
\hline
PATCH & /api/v1/tickets/\{id\}/status &
Cambia el estado de la solicitud & Técnico, Admin \\
\hline
POST  & /api/v1/tickets/\{id\}/assign &
Asigna la solicitud a un técnico & Admin \\
\hline
POST  & /api/v1/tickets/\{id\}/comments &
Agrega un comentario al caso & Todos \\
\hline
\end{longtable}

\section{Servicio de reportes}

\begin{longtable}{|p{1.8cm}|p{5.2cm}|p{3.8cm}|p{1.8cm}|}
\hline
\textbf{Método} & \textbf{Ruta} & \textbf{¿Qué hace?} &
\textbf{Acceso} \\
\hline
\endfirsthead
\hline
\textbf{Método} & \textbf{Ruta} & \textbf{¿Qué hace?} &
\textbf{Acceso} \\
\hline
\endhead
GET & /api/v1/reports/volume &
Cuántas solicitudes hubo en un período & Admin \\
\hline
GET & /api/v1/reports/technicians &
Rendimiento por técnico & Admin \\
\hline
GET & /api/v1/reports/status &
Solicitudes agrupadas por estado & Admin \\
\hline
\end{longtable}

\newpage

\chapter{Contenerización con Docker}

\section{¿Qué es Docker?}

Docker es una herramienta que permite empaquetar un servicio junto
con todo lo que necesita para funcionar: su código, sus librerías y
su configuración. Ese paquete se llama contenedor.

La ventaja es que un contenedor funciona exactamente igual en
cualquier computadora donde Docker esté instalado. Esto elimina el
problema clásico del desarrollo de software: ``en mi computadora
sí funciona''.

\section{Importancia de Docker}

El sistema tiene cinco microservicios, cada uno con tecnologías
distintas: algunos en Java, uno en Python, uno en Node.js. Sin Docker,
cada integrante del equipo tendría que instalar y configurar todas
esas tecnologías manualmente en su computadora, con el riesgo de que
las versiones no coincidan y algo falle.

Con Docker, cada servicio viene empaquetado y listo. Con un solo
comando se levanta todo el sistema completo en cualquier máquina,
sin instalaciones adicionales.

\section{Contenedores}

Cada uno de los cinco microservicios tendrá su propio contenedor.
Las bases de datos también corren en contenedores separados.
Todos se conectan a través de una red interna privada, de modo que
se pueden comunicar entre sí usando el nombre del servicio como
dirección, sin necesidad de IPs fijas.
\newpage

\begin{longtable}{|p{4.5cm}|p{3cm}|p{5cm}|}
\hline
\textbf{Contenedor} & \textbf{Puerto expuesto} & \textbf{Visibilidad} \\
\hline
\endfirsthead
\hline
\textbf{Contenedor} & \textbf{Puerto expuesto} & \textbf{Visibilidad} \\
\hline
\endhead
api-gateway          & 8080 & Accesible desde el exterior \\
\hline
auth-service         & 8001 & Solo en desarrollo \\
\hline
ticket-service       & 8002 & Solo en desarrollo \\
\hline
report-service       & 8005 & Solo en desarrollo \\
\hline
notification-service & ---  & Solo red interna \\
\hline
ai-service           & ---  & Solo red interna \\
\hline
\end{longtable}

En un entorno de producción real, solo el API Gateway estaría
expuesto hacia el exterior. Los demás serían invisibles desde fuera.


\newpage

\chapter{Conclusiones}

Esta segunda entrega convirtió la propuesta general de la Entrega 1
en un diseño técnico concreto. El equipo pasó de tener una idea de
alto nivel a tener decisiones específicas y documentadas sobre cómo
está organizado el sistema, cómo se comunican sus partes, cómo
gestiona sus datos y cómo se despliega.

Los resultados más importantes de esta etapa son:

\begin{enumerate}
  \item El sistema quedó dividido en cinco microservicios con
        responsabilidades claras y bases de datos independientes,
        eliminando el acoplamiento que existía en la aplicación
        original.

  \item Se definió un esquema de comunicación claro: REST para
        operaciones que necesitan respuesta inmediata y Kafka para
        eventos que pueden procesarse en segundo plano.

  \item El diseño de la base de datos original fue analizado y
        distribuido entre los servicios correspondientes, respetando
        la estructura ya existente del sistema.

  \item La estrategia de contenerización con Docker garantiza que
        cualquier integrante del equipo pueda levantar el sistema
        completo de forma reproducible, sin depender de
        configuraciones manuales.

  \item El contrato del API REST quedó definido antes de comenzar
        la implementación, lo que permitirá que el equipo trabaje
        en paralelo durante la Semana 7.
\end{enumerate}

\newpage

\section{Repositorio en GitHub}

El programa de esta práctica práctica se encuentra disponible en:

\url{https://github.com/CristhianP03/Proyecto-Grupal-App-Distribuidas-EquipoACC.git}


\section*{Rol de la Inteligencia Artificial}
Durante la elaboración de este informe se utilizó una herramienta de
inteligencia artificial como apoyo en la redacción. Su participación se limitó
a correcciones menores de texto y estructura, como mejorar la redacción de
algunos párrafos, reorganizar el orden de ciertas explicaciones y unificar
el estilo del documento.

\end{document}
